\documentclass[10pt]{extarticle}
% margins
\usepackage[letterpaper,top=0.4in, bottom=0.8in, left=0.9in, right=0.9in]{geometry}
% font
%\usepackage[cmintegrals,cmbraces]{newtxmath}
%\usepackage{ebgaramond-maths}
%\usepackage[T1]{fontenc}
% comment blocks
\usepackage{verbatim}
% no indent for new paragraphs and sections
\setlength{\parindent}{0in}
% \newsectionwidth{0pt}
% links and their colours
\usepackage[svgnames]{xcolor} % see https://www.latextemplates.com/svgnames-colors
\usepackage{hyperref}
\definecolor{accent}{HTML}{164480}
\hypersetup{colorlinks=true,urlcolor=accent}
% symbols
\usepackage{fontawesome}
% tables
\usepackage{xltabular}
% no page numbers
% \pagenumbering{gobble}
\pagestyle{plain}
% no widows
\usepackage{needspace}
% enumerate but backwards
\usepackage{etaremune}
% line spacing
\usepackage{setspace}

% dont stop when rendering
\nonstopmode
% set line spacing
\setstretch{1.1}

% left-center-right for header
\newcommand{\headerline}[3]{%
  \par\medskip\noindent
  \makebox[0pt][l]{#1}%
  \makebox[\textwidth][c]{#2}%
  \makebox[0pt][r]{#3}
  \par\medskip
}

% hrule with custom vertical spacing
\newcommand{\hrulesp}[1]{%
  \vspace{-1.5mm}
  \hrule height 0.15mm
  \vspace{#1}
}

\begin{document}

\headerline{%
  \faGlobe: \href{https://jshen.ca}{\texttt{jshen.ca}}}{%
  }{%
  \faGithub: \href{https://github.com/al-jshen}{\texttt{github.com/al-jshen}}%
}
\vspace{-5.5mm}
\headerline{%
  \faEnvelopeO: \href{mailto:jshenschool@gmail.com}{\texttt{jshenschool@gmail.com}}}{%
  \huge Jeff Shen}{%
\faLinkedin: \href{https://linkedin.com/in/jeff-shen}{\texttt{linkedin.com/in/jeff-shen}}}%
}

\needspace{5\baselineskip}
\vspace{-4mm}
\section*{Education}
\hrulesp{2mm}

\textbf{H.B.Sc., University of Toronto} \hfill 2018-2022 \\
Astronomy \& Astrophysics, Statistics, Mathematics. Dean's List Scholar.

\needspace{5\baselineskip}
\vspace{-4mm}
\section*{Publications}
\hrulesp{2mm}

\begin{etaremune}
\item \textbf{Jeff Shen}, Gwendolyn M. Eadie, Norman Murray, Dennis Zaritsky, Joshua S. Speagle, Yuan-Sen Ting, Charlie Conroy, Phillip A. Cargile, Benjamin D. Johnson, Rohan P. Naidu, Jiwon Jesse Han. The Mass of the Milky Way from the H3 Survey, submitted to ApJ.
\item \textbf{Jeff Shen}, Allison W. S. Man, Johannes Zabl, Zhi-Yu Zhang, Mikkel Stockmann, Gabriel Brammer, Katherine E. Whitaker, and Johan Richard. Molecular gas in a gravitationally lensed galaxy group at $z=2.9$. 2021, ApJ, 917, 79. \href{https://arxiv.org/abs/2105.11572}{2105.11572}.
\end{etaremune}

% \needspace{5\baselineskip}
\vspace{-7mm}
\section*{Research Experience}
\hrulesp{-2mm}

\begin{xltabular}{\textwidth}{lX}
  2021 Sep - Present & \textbf{Undergraduate Research Course} \\*
                     & \textit{Department of Astronomy and Astrophysics, University of Toronto} \\*
                     & Supervisor: \href{https://www.robertoabraham.com/}{Prof. Roberto Abraham} \\*
                     & Software and hardware for the Dragonfly Telephoto Array. \\
  2021 Jul - Present & \textbf{Independent Research} \\*
                     & Supervisors: \href{https://www.cita.utoronto.ca/~murray/}{Prof. Norman Murray} \\*
                     & Using Stan to fit dynamical models to geological data from thermal tides to estimate the length of day. \\
  2021 May - Present & \textbf{Independent Research} \\*
                      & Supervisors: \href{https://joshspeagle.github.io/}{Dr. Josh Speagle} \\*
                      & Probabilistic predictions of stellar ages with normalizing flows. Quantifying the contribution of stellar and Galactic information to age estimates. \\
  2021 May - 2021 Aug & \textbf{Undergraduate Summer Research Fellowship} \\*
  (Continued work     & \textit{Canadian Institute for Theoretical Astrophysics} \\*
  until present)      & Supervisors: \href{https://www.astro.utoronto.ca/~eadie/cover.html}{Prof. Gwendolyn Eadie}, \href{https://www.cita.utoronto.ca/~murray/}{Prof. Norman Murray}, \href{https://joshspeagle.github.io/}{Dr. Josh Speagle} \\*
                      & Continued work on the mass distribution of the Milky Way. Collaborated with members of the \href{http://h3survey.rc.fas.harvard.edu/}{H3 Survey} and submitted a first-author paper to ApJ. \\
  2020 Aug - 2021 May & \textbf{Undergraduate Research Course} \\*
                      & \textit{Canadian Institute for Theoretical Astrophysics} \\*
                      & Supervisors: \href{https://www.astro.utoronto.ca/~eadie/cover.html}{Prof. Gwendolyn Eadie}, \href{https://www.cita.utoronto.ca/~murray/}{Prof. Norman Murray} \\*
                      & Developed a hierarchical Bayesian model for modelling the mass of galactic halos using a fast and scalable sampling algorithm. \\
  2020 May - 2020 Aug & \textbf{Summer Undergraduate Research Program} \\*
  (Continued work  & \textit{Dunlap Institute for Astronomy and Astrophysics, University of Toronto} \\*
  until May 2021)  & Supervisor: \href{https://phas.ubc.ca/users/allison-man}{Prof. Allison Man} \\*
                   & Analyzed ALMA data to infer molecular gas masses in several gravitationally lensed high-redshift galaxies. Presented results in a first-author publication in ApJ.
\end{xltabular}


% \needspace{5\baselineskip}
\section*{Presentations}
\hrulesp{-2mm}

\begin{xltabular}{\textwidth}{lX}
  2021 August & \textbf{SDSS 2021 Collaboration Meeting, JHU/Online} \\*
           & Lightning talk: ``Predicting the ages of stars with machine learning.'' \\
  2021 May & \textbf{Stellar Stats Workshop, UofT/Online} \\*
           & Lightning talk: ``Estimating the mass distribution of the Milky Way with Bayesian multilevel models.'' \\
  2021 May & \textbf{Canadian Astronomical Society (CASCA) AGM, NRC Herzberg/Online} \\*
           & Poster presentation and lightning talk: ``Molecular gas in a gravitationally lensed galaxy group at z = 2.9.'' \\
  2021 Apr & \textbf{Vancouver Area Cosmology Meeting, UBC/SFU/TRIUMF/Online} \\*
           & Talk: ``Molecular gas in gravitationally lensed galaxies at z = 2.9.'' \\
  2020 Oct & \textbf{REQUIEM Galaxies Telecon, Online} \\*
           & Talk: ``Molecular gas in gravitationally lensed galaxies at z = 2.9.'' \\
  2020 Aug & \textbf{Summer Undergraduate Research Program Poster Session, UofT/Online} \\*
           & Poster presentation: ``Molecular gas in a gravitationally lensed galaxy group at z = 2.9.'' Awarded prize for one of top three posters. \\
\end{xltabular}


% \needspace{5\baselineskip}
\section*{Honours and Awards}
\hrulesp{-2mm}

\begin{xltabular}{\textwidth}{lX}
  2021 May - 2021 Aug & \textbf{Undergraduate Summer Research Fellowship, (\$8,396 CAD)} \\*
  2020 May - 2020 Aug & \textbf{Summer Undergraduate Research Program Award (\$9,500 CAD)} \\*
                      & Dunlap Institute for Astronomy and Astrophysics, University of Toronto
\end{xltabular}


% \needspace{5\baselineskip}
% \section*{Independent Coursework}

% \begin{xltabular}{\textwidth}{lX}
%     2020 & \textbf{Practical Time Series Analysis}, The State University of New York (Coursera) \\
%    & Stationarity, differencing, ACF, PACF, autocovariance. MA processes, AR processes, AR-MA duality. Backshift operator, difference equations, Yule-Walker matrix equations. AIC, SSE, model selection. ARMA and ARIMA models. Seasonality, SARIMA. Exponential smoothing, forecasting. \\
%     2019 & \textbf{Data-Driven Astronomy}, The University of Sydney (Coursera) \\
%     & Pulsars, black holes, exoplanets, stellar evolution, distances, redshift, galaxy classification. Image stacking, algorithm analysis, crossmatching, databases, regression, classification, machine learning. Implementation in Python with Astropy; SQL. \\
%     2019 & \textbf{Intro to Deep Learning with PyTorch}, Udacity \\
%     & Mathematics of neural networks. Data preparation, regularization techniques. Image classification with CNNs and OpenCV. Time series with RNNs. Implementation in Python with PyTorch. \\
%     2018 & \textbf{Intro to Machine Learning}, Udacity \\
%     & Classifiers (na\"ive Bayes, SVMs, ensemble methods), data processing, feature selection, regression, clustering, outlier detection, PCA, evaluation metrics. Implementation in Python with scikit-learn.\\
% \end{xltabular}

\needspace{5\baselineskip}
\section*{Selected Projects}
\hrulesp{2mm}

\begin{etaremune}
\item Comprehensive scientific computing package for Rust: linear algebra, statistics, numerical optimization, linear models, and more. \url{github.com/al-jshen/compute} 
\item Zero-dependency reverse-mode automatic differentiation package. \url{github.com/al-jshen/reverse} 
\item Fast and easy random number generation in Rust with Wyrand. \url{github.com/al-jshen/alea} 
% \item Orientation estimation on STM32 MCU (embedded system) with real-time visualizer. \url{github.com/al-jshen/rotvis}  
\item Parallelized ray tracer for realistic scene rendering. \url{github.com/al-jshen/traycer} 
\end{etaremune}

\needspace{7\baselineskip}
\vspace{-7mm}
\section*{Media}
\hrulesp{-2mm}

\begin{xltabular}{\textwidth}{lX}
  2021 Jul & \textbf{National Ratio Astronomy Organization (NRAO) eNews} \\
           & Digital newsletter feature: ``Molecular gas in high redshift galaxies''. \url{https://science.nrao.edu/enews/14.7/index.shtml#molecular_gas} \\
  2020 Jul & \textbf{SURP Student of the Week Feature} \\
           &  Interview: \url{https://www.dunlap.utoronto.ca/surp-sow/sow12020/} \\
\end{xltabular}


\needspace{5\baselineskip}
\section*{Outreach}
\hrulesp{-2mm}

\begin{xltabular}{\textwidth}{lX}
  2019 - Present & \textbf{Space Science Journalist (Current in Space segment)}, \href{https://thestarspot.ca/}{\textit{The Star Spot} podcast} \\
                 & Curate, write and record stories about the latest research and news in astronomy. Episodes 177 - present.
\end{xltabular}

\needspace{5\baselineskip}
\section*{Volunteering}
\hrulesp{-2mm}

\begin{xltabular}{\textwidth}{lX}
  2015 - 2018 & \textbf{Python Workshop Lead}, Sir Winston Churchill Secondary, Vancouver, B.C. \\
              & Coordinated and delivered workshops to prepare students for national programming contest.
\end{xltabular}

\needspace{5\baselineskip}
\section*{Technical Skills}
\hrulesp{2mm}

\textbf{Data}: Stan, Pyro, Pytorch, NumPy, SciPy, Pandas, Scikit-Learn, Matplotlib, SQL, R \\
\textbf{Other}: Rust, BLAS/LAPACK, Astropy, \LaTeX, regex, bash, git, Docker, AWS, Javascript, React


% \needspace{5\baselineskip}
% \section*{Relevant Coursework}

% \begin{xltabular}{\textwidth}{lX}
%     2020 Fall & Real Analysis \\
%     2020 Fall & Regression Analysis \\
%     2020 Fall & Practical Astronomy \\
%     2020 Winter & Galaxies and Cosmology \\
%     2020 Winter & Practical Physics I \\
%     2020 Winter & Electricity and Magnetism \\
%     2020 Winter & Thermodynamics and Statistical Mechanics \\
%     2019 - 2020 & Probability and Statistics I, II \\
%     2019 Fall & Stars and Planets \\
%     2019 Fall & Ordinary Differential Equations \\
%     2019 Summer & Multivariable Calculus \\
%     2019 Winter & Computer Science \\
%     2018 - 2019 & Calculus \\
%     2018 - 2019 & Linear Algebra I, II \\
%     2018 Fall & Statistical Reasoning and Data Science \\
% \end{xltabular}

% \needspace{5\baselineskip}
% \section*{Interests}

% \textbf{Languages}: English (native), French, Chinese, Norwegian (basic) \\
% \textbf{Other}: \mbox{*}nix customization, computer graphics, 3D modelling, running, baking, reading

\end{document}
